El SCDEE se ha concebido como un sistema completo con dos extremos
diferenciados pero acoplados. El \dfn{backend}, objeto de este trabajo,
gestiona la lógica de negocio, la persistencia de datos, los procesos
asíncronos de reconocimiento y notificación, la integración con el
almacenamiento de objetos y la seguridad. El \dfn{frontend}, objeto del
\acs{tfg} de Alejandro Soler Sichling, se centra en la experiencia de
usuario del corrector y del alumnado: visualización de páginas
escaneadas, anotación interactiva, gestión de calificaciones y consulta de
exámenes corregidos.

El contrato entre ambos extremos es la \ac{api} \ac{rest} versionada bajo
el prefijo \texttt{/api/v1/}, cuya especificación OpenAPI 3.0 se genera
automáticamente a partir del código del \dfn{backend} y constituye la
documentación de referencia para el desarrollo del \dfn{frontend}
(\cref{AP:OPENAPI}). Las decisiones técnicas que afectan a ambos extremos
se han discutido conjuntamente durante todo el desarrollo, asegurando una
arquitectura coherente.
